% Transcription — fonds Alexandre Grothendieck, shelfmark 15, batch 9 (pages 161-180).
% Established from the facsimile archives/batches/15/batch-09.pdf (a mirror of
% https://grothendieck.umontpellier.fr/15.pdf), per the skill
% transcribe-grothendieck. The batch file carries no cover sheet: PDF page N
% is page 160+N of the folder (the Montpellier footer prints « 161 » ... « 180 »),
% and the archivists' pencil numbers bottom-left agree leaf by leaf wherever
% they are legible (161-167, 169, 170, 172-178; 168, 171, 179 and 180 carry
% none that can be read, but sit in sequence).
%
% Pass: Opus 5.5 (claude-opus-5-5), 2026-09-23 — first pass, unchecked against the pages by a human.
%
% Eleven of the twenty pages are transcribed (162, 164, 167, 169, 172-178).
% Absent: pages 161, 163, 165, leaves of an unrelated typescript (a Bourbaki
% draft, « n° 413 », pp. 1-38, 1-37, 1-36, on point distributions and their
% convolution on a group manifold, Prop. 19-20; scanned upside down, nothing
% in his hand) — pages 162 and 164 are their versos; pages 166 and 168, blank;
% page 170, a leaf of another typescript (« Systèmes de racines (20) », Tribu
% n° 59, exercises on W(E7) and Sp(3,F2)), nothing in his hand; page 180, a
% computer listing of names and numbers, scanned upside down; page 171, an
% otherwise blank buff sheet inscribed top right in his hand « Calculs de
% cohomologie absolue des Motifs de Tate Z(n). » — a cover, whose words become
% the \section heading before page 172, per references/hand.md §5; page 179,
% an otherwise blank sheet carrying in pencil, centred, « Calculs de
% cohomologie absolue des motifs de Tate Z(n) » again — a second cover of the
% same words, not given a \page; a closing \note records it. Whether its
% pencil hand is his is not certain.
%
% What the batch is. Five runs, none paginated by him, in no evident order.
%   (1) Pages 162 and 164 (versos of the Bourbaki draft): Hodge
%   realisations of motives over a number field K. 162: « Si ξ̄ est réelle,
%   on a une structure de Hodge à involution. Comment se relient ces
%   structures de Hodge ?? », a diagram M̄_0 (motives over K̃ coming from K)
%   → G^0-Modules → Réseaux de Hodge by the fibre functors at two embeddings
%   ξ̃, ξ', then a long passage cancelled (horizontal strokes and two
%   diagonals), largely illegible, whose end is a first draft of 164. 164:
%   for V a G^0-module of weight p, an embedding ξ defines a K̄-subspace
%   H_ξ(V) ⊂ V_C with its filtration H_ξ(V)^{p,q}, making V a « réseau de
%   Hodge » whose Galois group is exactly G^0; then for V̄_ξ = V ⊗ K̄_ξ, the
%   H_ξ(V)^p ⊗_K K_ξ give a Hodge structure on (V, V_K) with group Im G^0,
%   and for ξ real, Frob_ξ ∈ G(Q) gives an involution of V, acting as the
%   identity on H_ξ(V) and compatible with the Hodge structure.
%   (2) Page 167, « Unités », on squared paper in blue ink, in a round and
%   careful hand quite unlike his elsewhere in the batch: for K a number field
%   in C, every unit has a real power iff K ⊂ R or K = K_0(α) with K_0
%   totally real and α² totally negative (a CM field); proof by comparing
%   unit ranks (Dirichlet), example Q(ζ).
%   (3) Page 169, headed « Foncteurs fibres motiviques, et groupes
%   fondamentaux motiviques. »: a programme « À examiner », items 0)-9)
%   (Weil abelian varieties attached to pure motives of odd weight,
%   Hodge-Riemann conditions; extensions of groups by motivic points; class
%   field theory; finite fields and motivic points with values in R; the role
%   of the motivic fundamental groups in comparing canonical lattices, Hodge
%   and ℓ-adic; L-functions and the functional equation; positive structures;
%   topology with coefficients in Z; Riemann hypothesis), with « OK » and
%   asterisks in the margin; a closing line on Hodge positivity by reduction
%   to divisors.
%   (4) Pages 172-177, after the cover of page 171, « Calculs de cohomologie
%   absolue des Motifs de Tate Z(n) »: 172 cohomology of π = Ẑ (k = F_q) with
%   coefficients Z_ℓ(n), Q_ℓ(n), then the absolute cohomology ℍ^i(F_q, Q(n)),
%   ℍ^i(F_q, Z(n)), and ℍ^i(Spec Z, Z(1)) = Z for i = 3, 0 otherwise; 173
%   the localisation sequence on X = Spec Z (i_x^! M = i_x^* M(-1)[-2]), for
%   M = Q(n): H^i(X,Q(n)) ≅ H^i(U,Q(n)) for n ≠ 1 (« est-ce vrai ?? »), and
%   for n = 1, H^i(X,Q(1)) = 0 except i = 3 « où on trouve Q ! »; 174
%   (cancelled by two diagonals) H^*(Ẑ, M), crossed homomorphisms φ(x^n),
%   Z^1, B^1, H^1(Z, M) = M/(1-f)M; 176 the same sequence for G_m on X̄ and
%   Ū, H^i_S(X, G_m) = H^{i-1}(S, Z), ending (Q/Z)^S → Q/Z; 177 H^i(Ū, G_m),
%   H^i(Ū, μ_{ℓ^n}), H^i(Ū, Z_ℓ(1)) and again H^i(X, Q(1)) = Q for i = 3 only.
%   (5) Pages 175 and 178: two leaves of a written-out redaction on the
%   Brauer group, numbered in §5 (5.1-5.6), in the order 178 then 175 by
%   their numbering though filed the other way: 178 (a framed passage on the
%   Leray spectral sequence of y → X cancelled by two diagonals), the exact
%   sequence (5.1) 0 → G_m → R^*_X → Div_X → 0 and H^i(X, Div_X) ≅
%   H^{i+1}(X, G_m) mod torsion for i ≥ 1, ending mid-sentence; 175
%   H^1(X, Div_X) = 0 when the strict local rings are factorial, Prop. 5.4
%   (H^2(X,G_m) → H^2(X,R^*_X) → Br(y) injective), Cor. 5.5 (Br(X) → Br(y)
%   injective, Br(X) ⊂ Br(K)), Th. 5.6 a) (ξ ∈ H^2(X,G_m) is in Br(X) iff
%   killed by a finite étale surjective f : X' → X), breaking off.
%
% Notation fixed for the batch. ξ, ξ̃, ξ' embeddings of K; K̄ and K̃ as he
% writes them (the bar and the tilde both occur); G^0 for his « G° »;
% \mathbb{H} for his double-barred H of absolute cohomology on page 172
% (plain H elsewhere, as written); \underset{\sim}{\mathbb{Q}},
% \underset{\sim}{\mathbb{Z}} for the Q, Z underlined with a wave on page
% 172; \mathfrak{U} for the looped U of page 174 (the open subgroup of π);
% \mathcal{D}\mathrm{iv}_X and \mathcal{H}^i_x for the sheaves he underlines
% (« Div » on 175 and 178, the local cohomology sheaf on 176); \mathbb{G}_m,
% \mu_{\ell^n}. Plain square brackets in the text are his; brown-ink
% additions (pp. 172, 175, 178) are not distinguished from the black. Page
% numbers in the text are written plainly (no \ref).
%
% Legibility. The formulas of 172-177 read almost entirely; the connective
% prose of 162 and 164, the cancelled passages of 162 and 178 and several
% words of 169 are partly \ill{}. Nothing on these leaves is dated.
%
% Cross-batch. Page 162 opens mid-thought (« Si ξ̄ est réelle, ... »), and
% may continue a page of batch 8; page 178 ends mid-sentence (« ... un
% isomorphisme de faisceaux étales : »), its continuation not in this batch;
% page 175 cites (5.3.) and 5.1 a), which are not in this batch either.
%
% The dating is the inventory's own for the folder, copied verbatim. Its
% group, [10-18] « Théorie des motifs », is dated [à partir de 1958-à partir
% de 1982].
\documentclass[11pt,a4paper]{article}
% Le préambule commun à toutes les transcriptions du fonds Grothendieck.
%
% Il tient en une idée : **l'incertitude se compose, elle ne se résout pas.**
% Une transcription qui lisse un mot douteux a détruit la seule chose pour
% laquelle on la faisait — un lecteur doit pouvoir savoir, à chaque endroit,
% ce qui était sur le feuillet et ce qui a été ajouté. Les six macros qui
% suivent sont donc l'appareil critique, et non de la décoration.
%
% Les mêmes commandes sont reconnues par `scripts/render.mjs`, qui produit la
% vue de lecture du site. Une macro ajoutée ici doit l'être aussi là-bas,
% faute de quoi le rendu échouera bruyamment — ce qui est voulu : un rendu
% silencieusement faux est pire qu'un rendu absent.

\ProvidesPackage{grothendieck}[2026/08/08 Transcriptions du fonds Grothendieck]

\usepackage{amsmath,amssymb,amsthm}
\usepackage{mathrsfs}
\usepackage{stmaryrd}
% Les diagrammes commutatifs sont partout dans ce fonds : c'est la langue
% même de la géométrie algébrique de Grothendieck, et une transcription qui
% les rendrait en prose ne transcrirait rien.
\usepackage{tikz-cd}
% Français par défaut — c'est la langue des feuillets — mais l'anglais est
% chargé aussi : la traduction et le résumé sont des documents anglais, et la
% typographie française (espaces avant les deux-points) y serait une faute.
% Ces documents basculent par \selectlanguage{english} en tête de corps.
\usepackage[english,french]{babel}
\usepackage{fontspec}
\usepackage{geometry}
\geometry{a4paper,margin=25mm,marginparwidth=28mm}
\usepackage{marginnote}
\usepackage{xcolor}
% ulem, not soul. `soul`'s \st and \ul are built on a character-by-character
% scan that XeLaTeX's font machinery breaks: the document fails with an
% undefined control sequence rather than degrading. ulem does the same two jobs
% and is Unicode-engine safe. `normalem` stops it hijacking \emph, which the
% transcriptions use for Grothendieck's own emphasis.
\usepackage[normalem]{ulem}
\usepackage{hyperref}
\hypersetup{colorlinks=true,linkcolor=black,urlcolor=black,pdfborder={0 0 0}}

% --- Métadonnées du lot -----------------------------------------------------
% Reprises telles quelles par le script de rendu, qui les affiche en tête de
% la vue de lecture. Elles fixent ce que le fichier prétend couvrir : sans
% elles, une transcription flotte sans point d'attache dans un fonds de seize
% mille feuillets.

\newcommand{\folder}[1]{\gdef\grothFolder{#1}}
\newcommand{\batch}[1]{\gdef\grothBatch{#1}}
\newcommand{\leaves}[2]{\gdef\grothFirst{#1}\gdef\grothLast{#2}}
\newcommand{\foldertitle}[1]{\gdef\grothFoldertitle{#1}}
\gdef\grothFolder{?}\gdef\grothBatch{?}\gdef\grothFirst{?}\gdef\grothLast{?}\gdef\grothFoldertitle{}

% --- Le résumé --------------------------------------------------------------
%
% Il ouvre la lecture modernisée, sous le titre « Résumé », et il est composé
% en petit. Ce n'est pas un ornement : c'est le seul passage du document qui
% ne soit pas des mathématiques, et un lecteur doit pouvoir le reconnaître
% avant de l'avoir lu — puis le sauter s'il connaît déjà le sujet, ou s'y
% arrêter s'il ne le connaît pas. Le filet qui le suit marque la reprise du
% corps.

\newenvironment{resume}
  {\small\setlength{\parskip}{0.45em}}
  {\par\medskip\hrule\medskip}

% --- Mots-clés ---------------------------------------------------------------
%
% En anglais, en clôture du résumé de la lecture modernisée : le vocabulaire
% moderne sous lequel chercher ce que les feuillets construisent. Le manifeste
% du site (`npm run manifest`) les extrait pour étiqueter la cote entière —
% cette ligne est la source unique des tags, il n'y a pas de fichier à côté.
\newcommand{\keywords}[1]{\par\smallskip\noindent
  {\footnotesize\textsc{Keywords} --- \textit{#1}}\par}

% --- L'appareil critique ----------------------------------------------------

% Le numéro de feuillet, en marge. C'est l'ancre qui permet de régler un
% désaccord sur une formule en regardant *un* feuillet plutôt qu'une chemise
% de six cents. Le site s'en sert aussi pour tourner les pages du fac-similé
% quand on fait défiler la transcription.
\newcommand{\leaf}[1]{%
  \marginnote{\footnotesize\color{gray!70}\textsf{#1}}%
  \label{leaf:#1}\ignorespaces}

% Illisible : on ne devine pas. Un mot inventé qui se lit comme les autres est
% le pire résultat possible.
\newcommand{\ill}{\textcolor{red!60!black}{[\,\ldots\,]}}

% Lecture douteuse : le mot est donné, et signalé comme incertain.
\newcommand{\uncertain}[1]{\uline{#1}}

% Ajout de l'éditeur — une lettre manquante, un mot sous-entendu.
\newcommand{\add}[1]{\textcolor{blue!50!black}{[#1]}}

% Biffé par Grothendieck lui-même. Ce qu'il a rayé fait partie du document :
% une rature dit souvent où la pensée a changé de direction.
\newcommand{\struck}[1]{\sout{#1}}

% Note du transcripteur, sur l'état matériel du feuillet.
\newcommand{\note}[1]{\par\smallskip{\small\color{gray!60!black}\textit{[#1]}}\par\smallskip}

% Note marginale de Grothendieck, distincte du corps du texte.
\newcommand{\marginal}[1]{\par\smallskip\noindent\hspace{1em}%
  {\small\color{gray!60!black}#1}\par\smallskip}

% --- Titre ------------------------------------------------------------------

\AtBeginDocument{%
  \begin{center}
    {\large\bfseries Fonds Alexandre Grothendieck --- cote n\textsuperscript{o}~\grothFolder}\\[2pt]
    {\itshape\grothFoldertitle}\\[4pt]
    {\small feuillets \grothFirst--\grothLast\ (lot \grothBatch)}\\[2pt]
    {\footnotesize\color{gray!60!black}Transcription établie d'après le fac-similé numérisé
     par l'Université de Montpellier.\\
     Les lectures incertaines sont soulignées, les illisibles notées [\,\ldots\,],
     les ajouts entre crochets.}
  \end{center}
  \vspace{1em}}

\endinput


\folder{15}
\batch{9}
\pages{161}{180}
\dating{1965-[vers 1977]}
\watermark{Édition de démonstration}
\foldertitle{Théorie arithmétique. Théorie de Galois des motifs (1965) : notes manuscites (s.d.), tapuscrit (s.d.).}

\begin{document}

\page{162}
Si $\bar{\xi}$ \uncertain{est} réelle, on a une structure de Hodge à
involution.

\note{le mot entre $\bar{\xi}$ et « réelle » est surchargé}

Comment se relient ces structures de Hodge ??

\begin{tikzcd}[column sep=small, row sep=small, nodes={font=\scriptsize}]
\overline{\mathcal{M}}_{0} \arrow[rr, "\text{fibre en } \tilde{\xi}"] \arrow[rr, bend right=25, "\text{fibre en } \xi'"'] & & G^{0}\text{-Modules} \arrow[rr, "\Phi_{\tilde{\xi}}"] \arrow[rr, bend right=25, no head, "\varphi"'] & & \text{Réseaux de Hodge}
\end{tikzcd}

$\overline{\mathcal{M}}_{0}$ motifs définis sur $\widetilde{K}$ provenant de
$K$.

\note{le diagramme est redessiné d'après l'esquisse : la légende de
$\overline{\mathcal{M}}_{0}$ est écrite sous le symbole, à gauche ; devant
« $G^{0}$-Modules », un premier mot (« Variét. » ?) est biffé. Sur la page,
les deux flèches de la seconde ligne partent de $\overline{\mathcal{M}}_{0}$
et de sous « $G^{0}$-Modules » ; la seconde, marquée $\varphi$, est un simple
trait sans pointe, et leurs buts ne sont pas écrits}

\struck{\ill{} La donnée de $\xi$ \ill{} de $K^{*}$ et $N_{K^{*}/\mathbb{Q}}$ ;}
\struck{NB Le seul \ill{} ainsi \ill{} $K^{*}$-modules \ill{} $K$. \ill{}
Si $E^{0} \neq 0$, \ill{} réel. Les \ill{} $K^{*}$ opère \ill{} $\mathbb{Z}$.
\ill{}}

\struck{$V$ un $G^{0}$-module,}
\struck{$V \otimes_{\mathbb{Q}} \mathbb{C}$ muni des \ill{} structure de
Hodge — $H^{V}_{\xi}$ (cf.\ plan des $K$)}
\struck{et d'un sous-$\widetilde{K}$-espace vectoriel $H^{V}_{\xi}$}

\note{toute la moitié inférieure du feuillet, à partir de « NB », est biffée
de traits horizontaux, encadrée, puis barrée de deux longues diagonales ; le
passage est en grande partie illisible. $K^{*}$ y est souligné deux fois. Les
trois dernières lignes sont un premier état de la définition reprise au
feuillet 164}

\page{164}
$V$ un $G^{0}$-module [de poids $p$ (N.B. le \uncertain{même} \ill{}
\uncertain{des poids}-$\mathbb{Z}$)], $V_{\mathbb{C}} = V \otimes_{\mathbb{Q}}
\mathbb{C}$.

Alors un plong.\ \struck{\ill{}} $\xi$ permet de définir
[\uncertain{coh.\ de De Rham}] un sous-$\overline{K}$-espace vectoriel
\[
  H_{\xi}(V) \subset V_{\mathbb{C}},
\]
filtré par des $H_{\xi}(V)^{p,q}$, tel que \struck{l'on} on ait une
\emph{\uncertain{réduction} des scalaires} à $K$, et que la filtration
$H_{\xi}(V)^{p,q}_{\mathbb{C}}$ fasse de $V$ un \emph{réseau de Hodge},
$G^{0}$ étant \struck{alors} exactement \emph{le groupe de Galois} de la
structure de Hodge. Si le plong.\ est réel, alors

\note{les crochets de la première ligne sont de sa main : « de poids $p$ »
et la parenthèse sont ajoutés au-dessus de la ligne ; « coh.\ de De Rham »
est ajouté au crayon, très pâle, au-dessus de « de définir » ; «
$\overline{K}$- » est ajouté dans l'interligne. « plong. » abrège sans doute
« plongement » ; le feuillet s'interrompt sur « alors »}

\struck{$\Psi$} $V$ un \uncertain{$E$}-module,

$\xi$ un plong.\ de $K$
\[
  \overline{V}_{\xi} = V \otimes_{\mathbb{Q}} \overline{K}_{\xi} \supset
  H_{\xi}(V)^{(p)}
\]

\begin{tikzcd}[column sep=small, row sep=small]
K_{\xi_{1}} & K_{\xi_{2}} & K_{\xi_{3}} \\
 & K \arrow[ul, no head] \arrow[u, no head] \arrow[ur, no head] & \\
 & \mathbb{Q} \arrow[u, no head] &
\end{tikzcd}

Les $H_{\xi}(V)^{p} \otimes_{K} K_{\xi}$ définissent sur $(V, V_{K})$ une
structure de Hodge, \struck{le} le groupe de Galois étant \struck{\ill{}}
\struck{contenu} $\mathrm{Im}\, G^{0}$.

\note{le petit diagramme des corps est tracé à gauche, en marge ; l'indice
du troisième $K$ est mal formé. Dans $H_{\xi}(V)^{(p)}$, un premier exposant
est biffé}

Si $\xi$ est réelle, on a une involution dans $V$ provenant de l'élém.\
\struck{\ill{}} $\mathrm{Frob}_{\xi}$ de $G(\mathbb{Q})$ et \ill{} \ill{}
\uncertain{prolonge} $\mathrm{Frob}_{\xi}$ en \uncertain{forme}
\emph{semi-linéaire} de $\overline{V}_{\xi}$, $\mathrm{Frob}_{\xi}$
\uncertain{induit} l'identité sur $H_{\xi}(V)$ [et par suite
$\mathrm{Frob}_{\xi}$ est compatible avec la structure de Hodge…]

\note{« semi-linéaire » est écrit « ½linéaire », selon son usage}

\section{Unités}

\note{titre en tête du feuillet 167. Ce feuillet, à l'encre bleue sur papier
quadrillé perforé, est d'une écriture ronde et appliquée, très différente de
la sienne dans le reste du lot ; il est transcrit comme les autres}

\page{167}
\emph{Soit $K$ un corps de nombres algébriques plongé dans $\mathbb{C}$. Les
conditions suivantes sont alors équivalentes} : a) \emph{pour toute unité
$\varepsilon$ de $K$, il y a $n \in \mathbb{N}$, $n > 0$ tel que
$\varepsilon^{n}$ soit réel} ; b) \struck{\ill{}} \emph{ou bien $K \subset
\mathbb{R}$ ou bien $K$ se met sous la forme $K_{0}(\alpha)$, avec $K_{0}$
totalement réel, $\alpha^{2} \in K_{0}$, $\alpha^{2}$ totalement négatif}.

Supposons $K \not\subset \mathbb{R}$ et soit $K' = K \cap \mathbb{R}$ ; $K$
est extension \struck{quadratique} [de degré $p \geq 2$] de $K'$. Soient
$r_{1}$ (resp.\ $r'_{1}$) le nombre de conjugués réels de $K$, $r_{2}$
(resp.\ $r'_{2}$) la moitié du nombre de conjugués imaginaires de $K'$ ; on
a $[K : K'] = 2$, d'où $r_{1} + 2r_{2} = p(r'_{1} + 2r'_{2})$. Si a) est
satisfaite, le groupe des unités de $K'$ a même rang que celui de $K$, d'où
\struck{\ill{}} $r_{1} + r_{2} = r'_{1} + r'_{2}$ ; il en résulte que $0
\geq r_{1} + 2r'_{2}$, d'où $r_{1} = r'_{2} = 0$ : $K'$ est totalement réel
et $K$ totalement imaginaire. Réciproquement, si b) est satisfaite, on a
$r_{1} = r'_{2} = 0$, $r_{2} = r'_{1}$, $r_{1} + r_{2} = r'_{1} + r'_{2}$,
et le groupe des unités de $K'$ a même rang que celui de $K$, d'où a).

\note{les soulignements de l'énoncé sont rendus en italique. « de degré $p
\geq 2$ » remplace dans l'interligne « quadratique » biffé, et le $p$ de
$p(r'_{1} + 2r'_{2})$ est repassé d'une encre plus foncée, sur un premier
chiffre ; « $[K : K'] = 2$ » est resté tel quel. La ligne « $0 \geq r_{1} +
2r'_{2}$ » est transcrite comme écrite}

\emph{Exemple}. La condition a) est satisfaite si $K = \mathbb{Q}(\zeta)$,
$\zeta^{m} = 1$, car, $K/\mathbb{Q}$ étant abélienne, le corps des nombres
réels de $K$ est totalement réel tandis que $K$ est totalement imaginaire.

\section{Foncteurs fibres motiviques, et groupes fondamentaux motiviques.}

\note{titre de sa main, souligné, en tête du feuillet 169, d'une encre plus
brune que la liste}

\page{169}
À examiner :
\begin{itemize}
  \item[0)] VA à la Weil [encadré] associées aux motifs purs de poids impair
  sur \ill{}

  [Conditions de Hodge-Riemann …]
  \item[OK * (1)] Extension de groupes définie par des pts motiviques.

  \ill{} gpe fondamental motivique.
  \item[(2)] Relations avec le corps de classes
  \item[OK (3)] Cas du \uncertain{spectre} d'un corps fini. Pts motiviques
  \ill{} à val.\ dans $\mathbb{R}$ ???
  \item[* 4)] Rôle des groupes fondamentaux motiviques pour la comparaison des
  réseaux canoniques

  [dans Coh.\ de Hodge : 2 réseaux transcendants pour deux places diff. ;
  réseau transcendant et réseau de Hodge

  dans Coh.\ $\ell$-adique : 2 réseaux transcendants]
  \item[5)] Relations [de 4)] avec mes \ill{} anciennes
  \item[6)] Fonction $L$ : équation fonctionnelle
  \item[OK 7)] Structures positives …
  \item[8)] Topologie à \uncertain{coeff.} sur $\mathbb{Z}$ …
  \item[9)] Hyp.\ de Riemann ??
\end{itemize}

\note{« OK » et les astérisques sont dans la marge gauche, au droit des
items ; les numéros (1), (2), (3) sont cerclés, et le 8) est écrit sur un
autre chiffre. « Conditions de Hodge-Riemann … » est ajouté entre les lignes
0) et (1), et la fin de la ligne 0) touche le bord déchiré du feuillet. Dans
l'item 4), les deux lignes entre crochets sont séparées de leur titre par un
trait vertical}

Positivité à la Hodge en se ramenant à des questions de \emph{diviseurs}
par récurrence, \ill{} \uncertain{variétés de drapeaux} ??

\section{Calculs de cohomologie absolue des Motifs de Tate $\mathbb{Z}(n)$.}

\note{titre de sa main sur un feuillet par ailleurs vierge (feuillet 171),
qui sert de couverture aux feuillets qui suivent}

\page{172}
\[
  k = \mathbb{F}_{q} \;\text{---}\; \bar{k}, \qquad \pi \simeq
  \hat{\mathbb{Z}} \ \text{engendré par } f = \mathrm{frob}^{q}
\]
$M$ un $\mathbb{Z}_{\ell}$-module de t.f.\ sur lequel $\pi$ opère
continûment.
\[
\begin{aligned}
  H^{0}(\pi, M) &= M^{\pi} = M^{(1-f)} = \operatorname{Ker}(1-f)_{M} \\
  H^{1}(\pi, M) &= M_{\pi} = M_{(1-f)} = \operatorname{Coker}(1-f)_{M} = M/(1-f)M \\
  H^{i}(\pi, M) &= 0 \quad \text{si } i \geq 2
\end{aligned}
\]

\[
  H^{0}(\pi, \mathbb{Z}_{\ell}(n)) = 0 \ \text{si } n \neq 0, \quad =
  \mathbb{Z}_{\ell} \ \text{si } n = 0
\]
\[
\begin{aligned}
  H^{1}(\pi, \mathbb{Z}_{\ell}(n)) &\underset{\text{non canoniquement}}{\simeq}
  \mathbb{Z}_{\ell}/(1-q^{n})\mathbb{Z}_{\ell} \\
  &= \begin{cases}
    \mathbb{Z}_{\ell} & \text{si } n = 0 \\
    0 & \text{si } \ell \nmid (1-q^{n}) \ \text{i.e. } n \not\equiv 0 \ (w_{\ell}(q)) \\
    \mathbb{Z}/\ell^{v_{\ell}(1-q^{n})} & \text{si } \ell \mid 1-q^{n},\ n \neq 0 \ \text{i.e. } n \equiv 0 \ (w_{\ell}(q))
  \end{cases}
\end{aligned}
\]
\marginal{$w_{\ell}(q) =$ ordre de $q$ mod $\ell$ dans
$(\mathbb{Z}/\ell\mathbb{Z})^{*}$}

\note{la note marginale, en haut à droite et d'encre brune, est reliée par
un trait aux conditions « i.e. … », ajoutées de la même encre ; sous la
première, un mot est biffé}

\[
\left\{
\begin{aligned}
  & H^{i}(\pi, \mathbb{Q}_{\ell}(n)) = 0 \quad \text{si } n \neq 0, \ \text{quel que soit } i \\
  & H^{0}(\pi, \mathbb{Q}_{\ell}) = \mathbb{Q}_{\ell}, \quad H^{1}(\pi, \mathbb{Q}_{\ell}) = \mathbb{Q}_{\ell}
\end{aligned}
\right.
\]

\struck{$H^{i}(\pi, \mathbb{Z}(n)$}

\[
\left\{
\begin{aligned}
  & \mathbb{H}^{i}(\mathbb{F}_{q}, \underset{\sim}{\mathbb{Q}}(n)) = 0 \quad \text{si } n \neq 0, \ \text{quel que soit } i \\
  & \mathbb{H}^{0}(\mathbb{F}_{q}, \underset{\sim}{\mathbb{Q}}) = \mathbb{Q}, \quad \mathbb{H}^{1}(\mathbb{F}_{q}, \underset{\sim}{\mathbb{Q}}) = \underset{\sim}{\mathbb{Q}}.
\end{aligned}
\right.
\]

\[
\left\{
\begin{aligned}
  & \mathbb{H}^{0}(\mathbb{F}_{q}, \underset{\sim}{\mathbb{Z}}) = \mathbb{Z} \ (\text{mod } p\text{-groupes ?}), \quad \mathbb{H}^{1}(\mathbb{F}_{q}, \underset{\sim}{\mathbb{Z}}) = \mathbb{Z} \ (\text{mod } p\text{-groupes ?}) \\
  & \mathbb{H}^{0}(\mathbb{F}_{q}, \underset{\sim}{\mathbb{Z}}(n)) = 0 \ (\text{id}), \quad \mathbb{H}^{1}(\mathbb{F}_{q}, \underset{\sim}{\mathbb{Z}}(n)) = \mathbb{Z}/(1-q^{n})\mathbb{Z} \ (\text{mod } p\text{-groupes ?})
\end{aligned}
\right.
\]

\note{dans la dernière accolade, le $\underset{\sim}{\mathbb{Z}}$ de la
première ligne est écrit sur un $\underset{\sim}{\mathbb{Q}}$ ; sous les
deux $\underset{\sim}{\mathbb{Z}}(n)$ de la seconde, un trait renvoie à « $n
\neq 0$ ». Les « (mod $p$-groupes ?) » et « (id) » sont d'encre brune. Un
trait horizontal traverse ensuite la page}

\[
\begin{aligned}
  \mathbb{H}^{i}(\operatorname{Spec} \mathbb{Z}, \mathbb{Z}(1)) &=
  \begin{cases} 0 & \text{si } i \neq 3 \\ \mathbb{Z} & \text{si } i = 3 \end{cases} \\
  \mathbb{H}^{i}(\operatorname{Spec} \mathbb{Z}, \mathbb{Z}) &=
  \begin{cases} 0 & \text{si } i \neq 0 \\ \mathbb{Z} & \text{si } i = 0 \end{cases}
\end{aligned}
\]

\[
  \mathbb{H}^{i}(\operatorname{Spec}(\mathbb{Z}), \mathbb{Q}(n))
  \xrightarrow{\ \sim\ } \mathbb{H}^{i}(U, \mathbb{Q}(n))
\]
si $U$ ouvert de Zariski de $\operatorname{Spec} \mathbb{Z}$, et $n \neq 0$
(tout $i$).

\note{ces trois lignes, à droite, sont d'encre brune}

\page{173}
$\operatorname{Spec} \mathbb{Z}$

\note{suit un croquis : un trait horizontal, un point marqué $x$}

\[
  i_{x}^{!}(M) = i_{x}^{*}(M)(-1)\,[-2]
\]
\[
  \cdots \to H^{i-2}(x, i_{x}^{*}(M)(-1)) \to H^{i}(X, M) \to H^{i}(U, M) \to
  H^{i-1}(x, i_{x}^{*}(M)(-1)) \to \cdots
\]

\note{devant $i_{x}^{*}(M)$, dans la première formule, un trait isolé}

$M = \underset{\sim}{\mathbb{Q}}(n)$

\note{le $\mathbb{Q}$ est écrit sur un $\mathbb{Z}$}

\struck{$H^{i}$} $i \neq 1, 2, 3$ :
\[
  H^{i}(X, M) \xrightarrow{\ \sim\ } H^{i}(U, M)
\]
\marginal{\uncertain{pour des} \ill{}, car si $i \geq 4$, les deux membres
\uncertain{devraient être} $= 0$}
\[
\begin{aligned}
  0 &\to H^{1}(X, M) \to H^{1}(U, M) \to H^{0}(S, i^{*}(M)(-1)) \to H^{2}(X, M) \\
  &\to H^{2}(U, M) \to H^{1}(S, i^{*}(M)(-1)) \to H^{3}(X, M) \\
  &\to \underset{0\,?}{H^{3}(U, M)} \to 0
\end{aligned}
\]

$M =$ \struck{$\mathbb{Z}$} $\mathbb{Q}_{\ell}(n)$, $i^{*}(M)(-1) =
\mathbb{Q}(n-1)_{S}$

\note{l'indice $\ell$ de $\mathbb{Q}_{\ell}$ est peu sûr ; à droite du
premier membre, un signe biffé}

\emph{Si $n \neq 1$}
\[
  H^{i}(X, \mathbb{Q}(n)) \xrightarrow{\ \sim\ } H^{i}(U, \mathbb{Q}(n))
  \qquad \text{tout } i \quad [\text{est-ce vrai ??}]
\]

Si $n = 1$
\[
\begin{aligned}
  0 &\to H^{1}(X, \mathbb{Q}(1)) \to H^{1}(U, \mathbb{Q}(1)) \to \mathbb{Q}^{S} \to H^{2}(X, \mathbb{Q}(1)) \to \\
  &\to H^{2}(U, \mathbb{Q}(1)) \to \mathbb{Q}^{S} \to H^{3}(X, \mathbb{Q}(1)) \to 0
\end{aligned}
\]
$H^{i}(X, \mathbb{Q}(1)) = 0$ sauf pour $i = 3$, où on trouve $\mathbb{Q}$ !

\page{174}
\note{le feuillet entier est barré de deux longues diagonales}

\[
  k = \mathbb{F}_{q} \;\text{---}\; \bar{k}, \qquad \pi = \hat{\mathbb{Z}}
  \ \text{engendré par } \mathrm{frob}^{q}_{q}
\]
\[
\begin{aligned}
  H^{0}(\pi, M) &= M^{\pi} \\
  H^{1}(\pi, M) &= \varprojlim_{n} H^{1}(\pi, M/\ell^{n}M) =
\end{aligned}
\]
a) si $\pi$ opère trivialement, c'est $\mathrm{Hom}(\pi, M)$

b) si $M$ est fini, donc $\pi$ opère à travers $\pi/\mathfrak{U} \simeq
\mathbb{Z}/n\mathbb{Z}$
\[
  0 \to \underset{\substack{\parallel \\ M}}{H^{1}(\pi/\mathfrak{U}, M^{\mathfrak{U}})}
  \to H^{1}(\pi, M) \to
  \underset{\substack{\parallel \\ \mathrm{Hom}(\mathfrak{U}, M^{\pi})}}{H^{0}(\pi/\mathfrak{U}, \mathrm{Hom}(\mathfrak{U}, M))}
  \to \underset{\substack{\parallel \\ M}}{H^{2}(\pi/\mathfrak{U}, M^{\mathfrak{U}})}
\]
\struck{$M_{\pi}$} $\dfrac{\operatorname{Ker} N}{\operatorname{Im}(1-f)}$

\note{la dernière flèche de la suite descend vers $H^{2}$, écrit plus bas
à droite ; « $H^{1}(\pi, M)$ » est aussi écrit seul, à gauche, en marge}

\[
\begin{aligned}
  \varphi(xy) &= \varphi(x) + x\varphi(y) \\
  \varphi(x) &= xa - a \\
  \varphi(xy) &= xya - a = xa - a + x(ya - a)
\end{aligned}
\]
$N \quad 1 - f$, $\quad N = 1 + f + \cdots + f^{n-1}$
\[
\begin{aligned}
  \varphi(x^{2}) &= \varphi(x) + x\varphi(x) \\
  \varphi(x^{3}) &= \varphi(x) + x\varphi(x) + x^{2}\varphi(x) \\
  \varphi(x^{n}) &= (1 + x + \cdots + x^{n-1})\varphi(x) \\
  \text{or } \varphi(x^{-1}x) &= \varphi(x^{-1}) + x^{-1}\varphi(x) \\
  \varphi(x^{-1}) &= -x^{-1}\varphi(x) \\
  \varphi(x^{-n}) &= -(1 + x^{-1} + \cdots + x^{-(n-1)})x^{-1}\varphi(x)
  = -\frac{1 - x^{-n}}{1 - x^{-1}}\, x^{-1}\varphi(x) = \frac{1 - x^{-n}}{1 - x}
\end{aligned}
\]
\[
  \varphi(f^{n}) =
  \begin{cases}
    (1 + x + \cdots + x^{n-1})\varphi(x) = \dfrac{1 - x^{n}}{1 - x}\,\varphi(x) & \text{si } n \geq 1 \\
    0 & \text{si } n = 0
  \end{cases}
\]
\[
  \varphi(f^{n}) = 1 + x + \cdots + x^{n-1} \qquad \frac{1 - f^{n}}{1 - f}\,\varphi(f)
\]
\[
\begin{aligned}
  Z^{1}(\mathbb{Z}, M) &\simeq M \\
  B^{1}(\mathbb{Z}, M) &\simeq (1-f)M \\
  H^{1}(\mathbb{Z}, M) &\simeq M/(1-f)M
\end{aligned}
\]
si $n \equiv 0 \ (N)$, $\ 1 + f + \cdots + f^{n-1}$

\note{calculs épars, dans le désordre de la page : la colonne des
$\varphi(x^{n})$ est à gauche, sous la première diagonale ; quelques
égalités intermédiaires sont biffées (dont un premier $\varphi(x^{n}) =
\varphi(x) + \cdots$) ; la troisième ligne de l'accolade de $\varphi(f^{n})$
est commencée et laissée. La dernière égalité de la colonne s'arrête au
quotient, sans $\varphi(x)$}

\page{175}
Lorsque \struck{$X$ est} \struck{\ill{} on est dans les conditions de} les
anneaux locaux stricts de $X$ sont factoriels, on tire aisément de (5.3.)
\[
  H^{1}(X, \mathcal{D}\mathrm{iv}_{X}) = 0
\]
d'où \struck{la proposition}

\emph{Proposition 5.4.} Supposons que [l'on soit dans les conditions de 5.1.
a)], [$X$ étant irréductible de pt gén.\ $y$.] Alors les homomorphismes
canoniques
\[
  H^{2}(X, \mathbb{G}_{m,X}) \to H^{2}(X, R^{*}_{X}) \to H^{2}(y,
  \mathbb{G}_{m,y}) = \mathrm{Br}(y)
\]
sont injectifs.

\note{les deux crochets sont des ajouts de sa main dans l'interligne,
au-dessus de « Supposons que » et de « Alors les homomorphismes » ; « ($X$
étant irréductible …) » est d'abord écrit plus bas, puis remonté}

Comme $\mathrm{Br}(X)$ \uncertain{$\subset$} $H^{2}(X,
\mathbb{G}_{m,X})$, il s'ensuit

\emph{Corollaire 5.5.} Sous les mêmes conditions,
\[
  \mathrm{Br}(X) \to \mathrm{Br}(y)
\]
est injectif.

Ainsi, $\mathrm{Br}(X)$ s'identifie [\ill{}] à un sous-groupe du groupe de
Brauer $\mathrm{Br}(K)$ du corps des fonctions $K = k(y)$ de $X$.

\emph{Théorème 5.6.} [Soit $X$ un préschéma, $\xi \in H^{2}(X,
\mathbb{G}_{m})$] a) Si $X$ est un \uncertain{schéma local}, et
\struck{$\xi \in H^{2}(X, \mathbb{G}_{m,X})$}, pour que $\xi$
\struck{provienne d'un élément de} [soit dans $\mathrm{Br}(X)$], il faut et il
suffit qu'il existe un \struck{extension} morphisme fini étale surjectif
(galoisien si on veut) $f : X' \to X$, \struck{tel que $\xi | X'$ soit} qui
\struck{splitte} annule $\xi$.

\note{les ajouts de la fin, depuis « morphisme », sont d'encre brune ; le
feuillet s'arrête sur cette phrase}

\page{176}
\struck{$\mu_{n} \to \mathbb{G}_{m} \xrightarrow{n} \mathbb{G}_{m}$}

\struck{$\mu$}

\struck{$H^{3}(X, \mathbb{G}_{m}) \to H^{3}$}

\struck{$\mathcal{H}^{i}_{x}($}
\[
  H^{i}_{S}(X, \mathbb{G}_{m}) \to H^{i}(X, \mathbb{G}_{m}) \to H^{i}(U,
  \mathbb{G}_{m}) \to H^{i+1}_{S}(X, \mathbb{G}_{m}) \to
\]
\[
  \mathcal{H}^{i}_{x}(\mathbb{G}_{m}) =
  \begin{cases}
    0 & \text{si } i = 0 \\
    \mathbb{Z}_{x} & \text{si } i = 1 \\
    0 & \text{si } i \geq 2
  \end{cases}
\]
\[
  H^{i}_{S}(X, \mathbb{G}_{m}) = H^{i-1}(S, \mathbb{Z}) =
  \begin{cases}
    0 & \text{si } i = 0 \\
    \mathbb{Z}^{S} & \text{si } i = 1 \\
    0 & \text{si } i = 2 \\
    H^{1}(S, \mathbb{Q}/\mathbb{Z}) & \text{si } i = 3 \\
    0 & \text{si } i \geq 4
  \end{cases}
\]
\[
  0 \to \mathbb{Z} \to \mathbb{Q} \to \mathbb{Q}/\mathbb{Z} \to 0
\]

\note{les lignes $i = 1, 2, 3$ de la dernière accolade sont réunies à
droite par un crochet ; à gauche de la suite suivante, isolé, « $i = 1, 3$ ;
$i = 0, 2$ »}

Si $i \neq 0, 1, 2, 3$, $H^{i}(\overline{X}, \mathbb{G}_{m})
\xrightarrow{\ \sim\ } H^{i}(U, \mathbb{G}_{m})$ pour des \ill{}, car les deux
membres sont nuls.
\[
\begin{aligned}
  0 &\to \underset{\substack{\parallel \\ \mathbb{Z}^{*} = \{+1, -1\}}}{H^{0}(\overline{X}, \mathbb{G}_{m})}
  \to H^{0}(\overline{U}, \mathbb{G}_{m}) \to \mathbb{Z}^{S}
  \to \underset{\substack{\parallel \\ 0}}{H^{1}(\overline{X}, \mathbb{G}_{m})}
  \to \underset{\substack{\parallel \\ 0\,!}}{H^{1}(\overline{U}, \mathbb{G}_{m})}
  \to \underset{\substack{\parallel \\ 0}}{H^{2}_{S}(\overline{X}, \mathbb{G}_{m})} \\
  0 &\to \underset{\substack{\parallel \\ 0}}{H^{2}(\overline{X}, \mathbb{G}_{m})}
  \to H^{2}(\overline{U}, \mathbb{G}_{m})
  \to \underset{\substack{\parallel \\ (\mathbb{Q}/\mathbb{Z})^{S}}}{H^{1}(S, \mathbb{Q}/\mathbb{Z})}
  \to \underset{\substack{\parallel \\ \mathbb{Q}/\mathbb{Z}}}{H^{3}(\overline{X}, \mathbb{G}_{m})}
  \to \underset{\substack{\parallel \\ 0\,?}}{H^{3}(\overline{U}, \mathbb{G}_{m})} \to 0
\end{aligned}
\]

\note{sous la seconde suite, une flèche double relie $(\mathbb{Q}/\mathbb{Z})^{S}$
à $\mathbb{Q}/\mathbb{Z}$, marquée d'un mot au-dessus (« somme » ?) et d'un
autre au-dessous (« surject. » ?), tous deux peu lisibles. La barre de
$\overline{X}$ dans « Si $i \neq 0, 1, 2, 3$ » est incertaine}

\page{177}
\[
\begin{aligned}
  H^{0}(\overline{U}, \mathbb{G}_{m}) &= \mathbb{Z}/2 + \mathbb{Z}^{S} \\
  H^{1}(\overline{U}, \mathbb{G}_{m}) &= 0 \\
  H^{2}(\overline{U}, \mathbb{G}_{m}) &= \operatorname{Ker}\bigl((\mathbb{Q}/\mathbb{Z})^{S} \to \mathbb{Q}/\mathbb{Z}\bigr) \\
  H^{i}(\overline{U}, \mathbb{G}_{m}) &= 0 \quad \text{si } i \geq 3
\end{aligned}
\]

\note{la première ligne est d'encre brune}

\[
  0 \to \mu_{\ell^{n}} \to \mathbb{G}_{m} \xrightarrow{\ell^{n}}
  \mathbb{G}_{m} \to 0
\]
d'où
\[
\begin{aligned}
  H^{1}(\overline{U}, \mu_{\ell^{n}}) &\simeq (\mathbb{Z}/\ell^{n}\mathbb{Z})^{S} \\
  H^{2}(\overline{U}, \mu_{\ell^{n}}) &\simeq \operatorname{Ker}\bigl((\mathbb{Z}/\ell^{n}\mathbb{Z})^{S} \to (\mathbb{Z}/\ell^{n}\mathbb{Z})\bigr) \\
  H^{i}(\overline{U}, \mu_{\ell^{n}}) &= 0 \quad \text{si } i \geq 3
\end{aligned}
\]
\[
  H^{0}(\overline{U}, \mathbb{Z}_{\ell}) \simeq \mathbb{Z}_{\ell}
\]
\[
\left\{
\begin{aligned}
  H^{1}(U, \mathbb{Z}_{\ell}(1)) &\simeq \mathbb{Z}_{\ell}^{S} \\
  H^{2}(\overline{U}, \mathbb{Z}_{\ell}(1)) &\simeq \operatorname{Ker}(\mathbb{Z}_{\ell}^{S} \to \mathbb{Z}_{\ell})
\end{aligned}
\right.
\]
\marginal{petite modification pour $\ell = 2$ \uncertain{en dim.\ 0}}
\[
\left\{
\begin{aligned}
  H^{0}(X, \mathbb{Q}(1)) &= 0 \\
  H^{1}(X, \mathbb{Q}(1)) &= 0 \\
  H^{2}(X, \mathbb{Q}(1)) &= 0 \\
  H^{3}(X, \mathbb{Q}(1)) &= \mathbb{Q}
\end{aligned}
\right.
\]

\page{178}
\struck{On peut supposer $X$ irréductible de pt générique $y$. L'inclusion
$y \to X$ définit une suite spectrale (pour la cohomologie étale)}
\struck{$H^{*}(y, \mathbb{G}_{m}) \Longleftarrow E_{2}^{pq} = H^{p}(X,
R^{q}i_{*}(\mathbb{G}_{m}))$}
\struck{Or \ill{} on connaît les $H^{q}(y, \mathbb{G}_{m})$ et les
fibres de $R^{q}i_{*}(\mathbb{G}_{m})$, qui s'explicitent par de la
cohomologie galoisienne, et sont des groupes de torsion pour $q
> 0$, on trouve que \ill{} les groupes [sont un groupe de torsion
pour $p > 0$,]}
\struck{(*) $H^{p}(X, R^{*}_{X})$, où $R^{*}_{X} = R^{0}i_{*}(\mathbb{G}_{m})$} \struck{où $R^{*}_{X} =
i_{*}(\mathbb{G}_{m})$ est le faisceau [\ill{} cités] des fonctions
rationnelles [non nulles] sur $X$ \ill{}, si $X$ est
« sans idéaux premiers immergés »,}

\note{tout ce début est encadré d'un trait brun et barré de deux longues
diagonales ; les ajouts entre crochets y sont d'encre brune}

D'autre part, on a la suite exacte
\[
  (5.1.) \qquad 0 \to \mathbb{G}_{m,X} \to R^{*}_{X} \to
  \mathcal{D}\mathrm{iv}_{X} \to 0
\]
où $\mathcal{D}\mathrm{iv}_{X}$ est le faisceau (pour la top.\ [étale]) des
diviseurs [\uncertain{positifs}] (de Cartier) sur $X$, ce qui
\struck{\ill{}} [implique], compte tenu de (*), que
\[
  H^{i}(X, \mathcal{D}\mathrm{iv}_{X}) \xrightarrow{\ \partial\ } H^{i+1}(X,
  \mathbb{G}_{m,X})
\]
est bijectif mod torsion \struck{\ill{}} si $i \geq 1$. \struck{\ill{}}
\struck{\ill{} ne sera pas utilisé} L'hypothèse de factorialité d'autre part
implique \struck{l'isomorphisme} qu'\uncertain{on a un} isomorphisme de
faisceaux étales :

\marginal{\uncertain{Démonstration du corollaire} 5.4.}

\note{« $\mathcal{D}\mathrm{iv}$ » est souligné, pour un faisceau ; les
ajouts entre crochets et la fin du feuillet, depuis « L'hypothèse », sont
d'encre brune. La note marginale, à gauche en biais, est d'encre brune, deux
fois soulignée : le feuillet 175 appelle 5.4 une proposition. Le feuillet
s'arrête sur les deux-points ; la suite n'est pas dans ce lot}

\note{le feuillet 179, par ailleurs vierge, porte seul, au crayon, « Calculs
de cohomologie absolue des motifs de Tate $\mathbb{Z}(n)$ » : les mots de la
couverture du feuillet 171, repris}

\end{document}
